\documentclass[11pt,a4paper]{article}
\usepackage[utf8]{inputenc}
\usepackage{amsmath,amssymb,amsthm}
\usepackage{mathtools}
\usepackage{bm}
\usepackage{geometry}
\usepackage{hyperref}
\usepackage{cleveref}

\geometry{margin=1in}

\theoremstyle{definition}
\newtheorem{definition}{Definition}[section]
\newtheorem{theorem}{Theorem}[section]
\newtheorem{lemma}{Lemma}[section]
\newtheorem{proposition}{Proposition}[section]
\newtheorem{remark}{Remark}[section]

\title{Mathematical Proofs for Spatial Cramér-von Mises Test\\Formalization in Lean 4}
\author{Spatial-CvM Formalization Project}
\date{\today}

\begin{document}

\maketitle

\begin{abstract}
This document presents the mathematical proofs underlying the Spatial Cramér-von Mises (CvM) test formalization in Lean 4. We cover the asymptotic theory for spatially dependent data, including kernel smoothing, empirical processes, weak convergence, and the functional delta method. The proofs are designed to be compatible with the Mathlib library infrastructure.
\end{abstract}

\tableofcontents

\section{Introduction}

The Spatial Cramér-von Mises test assesses homogeneity of spatially dependent random fields. The formalization consists of three main theorems:
\begin{enumerate}
    \item \textbf{Theorem 1}: Weak convergence of the kernel-smoothed empirical process to a Gaussian process
    \item \textbf{Theorem 2}: Joint convergence of multiple spatial locations and the limiting chi-square distribution
    \item \textbf{Theorem 3}: Functional delta method for the test statistic
\end{enumerate}

\section{Preliminaries}

\subsection{Location Space}

\begin{definition}[Location Space]
Let $\mathcal{L} = \mathbb{R}^2$ be the location space with the Euclidean metric $d(s_1, s_2) = \|s_1 - s_2\|_2$.
\end{definition}

\subsection{Kernels}

\begin{definition}[Kernel Function]
A function $K: \mathbb{R} \to \mathbb{R}$ is a \emph{kernel} if it satisfies:
\begin{enumerate}
    \item \textbf{Symmetry}: $K(x) = K(-x)$ for all $x \in \mathbb{R}$
    \item \textbf{Boundedness}: There exists $B > 0$ such that $|K(x)| \leq B$ for all $x$
    \item \textbf{Compact Support}: There exists $r > 0$ such that $K(x) = 0$ for $|x| > r$
    \item \textbf{Lipschitz}: There exists $L > 0$ such that $|K(x) - K(y)| \leq L|x - y|$
    \item \textbf{Measurability}: $K$ is Borel measurable
\end{enumerate}
\end{definition}

\begin{definition}[Scaled Kernel]
For bandwidth $h > 0$, the scaled kernel is defined as:
\begin{equation}
K_h(x) = \frac{1}{h^2} K\left(\frac{x}{h}\right)
\end{equation}
\end{definition}

\section{Lemma 1: Asymptotic Covariance Structure}

\subsection{The Gamma Kernel}

\begin{definition}[Gamma Kernel]
For a kernel $K$, bandwidth $h > 0$, and lag $u \in \mathbb{R}$, the gamma kernel is:
\begin{equation}
\gamma_{K,h}(u) = \int_{\mathbb{R}} K_h(v) \cdot K_h(v - u) \, dv
\end{equation}
This is the convolution of the scaled kernel with itself.
\end{definition}

\begin{theorem}[Gamma Kernel at Zero]
For any kernel $K$ and $h > 0$:
\begin{equation}
\gamma_{K,h}(0) = \int_{\mathbb{R}} K_h(v)^2 \, dv > 0
\end{equation}
\end{theorem}

\begin{proof}
Since $K$ is not identically zero (otherwise it would not be useful as a kernel), there exists a set of positive measure where $K \neq 0$. By boundedness and compact support, $K_h \in L^2(\mathbb{R})$. The integral of a non-negative, not-identically-zero function over a set of positive measure is strictly positive.
\end{proof}

\subsection{The Gamma Operator}

\begin{definition}[Gamma Operator]
The covariance operator for the empirical process is:
\begin{equation}
\Gamma_{K,h}(s_1, s_2) = \int_{\mathbb{R}} K_h(u - s_1) \cdot K_h(u - s_2) \, du
\end{equation}
\end{definition}

\begin{theorem}[Symmetry of Gamma Operator]
For all $s_1, s_2 \in \mathbb{R}$:
\begin{equation}
\Gamma_{K,h}(s_1, s_2) = \Gamma_{K,h}(s_2, s_1)
\end{equation}
\end{theorem}

\begin{proof}
By commutativity of multiplication in the integrand:
\begin{align}
\Gamma_{K,h}(s_1, s_2) &= \int_{\mathbb{R}} K_h(u - s_1) \cdot K_h(u - s_2) \, du \\
&= \int_{\mathbb{R}} K_h(u - s_2) \cdot K_h(u - s_1) \, du \\
&= \Gamma_{K,h}(s_2, s_1)
\end{align}
\end{proof}

\subsection{Stationarity}

\begin{definition}[Stationary Random Field]
A spatial random field $\{X(s) : s \in \mathcal{L}\}$ is \emph{stationary} if for all shifts $\tau \in \mathcal{L}$:
\begin{equation}
X(s + \tau) \stackrel{d}{=} X(s)
\end{equation}
where $\stackrel{d}{=}$ denotes equality in distribution.
\end{definition}

\begin{theorem}[Stationary Substitution]
Let $f: \mathcal{L} \to \mathbb{R}$ be a function and $s_0 \in \mathcal{L}$, $h > 0$ parameters. Define:
\begin{equation}
g(u) = f(s_0 + h \cdot u)
\end{equation}
Then for all $s \in \mathcal{L}$:
\begin{equation}
f(s) = g\left(\frac{s - s_0}{h}\right)
\end{equation}
\end{theorem}

\begin{proof}
Direct computation:
\begin{equation}
g\left(\frac{s - s_0}{h}\right) = f\left(s_0 + h \cdot \frac{s - s_0}{h}\right) = f(s)
\end{equation}
\end{proof}

\section{Mixing and Dependence}

\subsection{Alpha-Mixing}

\begin{definition}[$\alpha$-Mixing Coefficient]
For a stationary random field $X$, the $\alpha$-mixing coefficient at distance $d$ is:
\begin{equation}
\alpha(d) = \sup \{|P(A \cap B) - P(A)P(B)| : A \in \mathcal{F}_{-\infty}^0, B \in \mathcal{F}_d^{\infty}\}
\end{equation}
where $\mathcal{F}_a^b$ is the $\sigma$-algebra generated by $\{X(s) : a \leq \|s\| \leq b\}$.
\end{definition}

\begin{definition}[Alpha-Mixing Condition]
A random field satisfies the \emph{alpha-mixing condition} with coefficient function $\alpha: \mathbb{R}_+ \to \mathbb{R}_+$ if:
\begin{enumerate}
    \item $\alpha(d) \geq 0$ for all $d \geq 0$
    \item $\sum_{d=1}^{\infty} \alpha(d) < \infty$ (summability)
\end{enumerate}
\end{definition}

\subsection{Davydov's Inequality}

\begin{theorem}[Davydov's Inequality]
Let $X$ and $Y$ be random variables with $\|X\|_p < \infty$ and $\|Y\|_q < \infty$ where $p, q \geq 2$ and $1/p + 1/q + 1/r = 1$. Then:
\begin{equation}
|\text{Cov}(X, Y)| \leq 4 \alpha(d)^{1/r} \|X\|_p \|Y\|_q
\end{equation}
where $d$ is the distance separating the $\sigma$-algebras to which $X$ and $Y$ are adapted.
\end{theorem}

\begin{proposition}[Non-negativity of Mixing Factor]
For $\alpha$-mixing coefficient $\alpha(d)$ and $q \geq 2$, the Davydov factor satisfies:
\begin{equation}
4 \cdot \alpha(d)^{1 - 2/q} \geq 0
\end{equation}
\end{proposition}

\begin{proof}
Since $\alpha(d) \geq 0$ by definition of mixing coefficients, and $4 > 0$, the product is non-negative. Real powers of non-negative numbers preserve non-negativity.
\end{proof}

\section{Theorem 1: Weak Convergence}

\subsection{The Empirical Process}

\begin{definition}[True Process]
For kernel $K$, bandwidth $h > 0$, true CDF $F$, and evaluation point $t$:
\begin{equation}
\tilde{F}_{K,h}(t) = \int_{\mathbb{R}} K_h(u - t) \cdot F(u) \, du
\end{equation}
\end{definition}

\begin{definition}[Empirical Process]
Given $n$ samples $\{X_i\}_{i=1}^n$ with empirical CDF $\hat{F}_n(u) = \frac{1}{n}\sum_{i=1}^n \mathbf{1}(X_i \leq u)$, the empirical process is:
\begin{equation}
\hat{Z}_n(t) = \sqrt{n} \int_{\mathbb{R}} K_h(t - u) \, d(\hat{F}_n(u) - F(u))
\end{equation}
\end{definition}

\subsection{Supremum Norm}

\begin{definition}[Supremum Norm]
For $f: \mathbb{R} \to \mathbb{R}$ and domain $D \subseteq \mathbb{R}$:
\begin{equation}
\|f\|_{\infty, D} = \sup_{x \in D} |f(x)|
\end{equation}
\end{definition}

\begin{proposition}[Supremum Norm Bound]
If $|f(x)| \leq B$ for all $x \in D$ and $D$ is nonempty, then:
\begin{equation}
\|f\|_{\infty, D} \leq B
\end{equation}
\end{proposition}

\subsection{Tightness}

\begin{definition}[Tightness]
A sequence of stochastic processes $\{X_n\}_{n=1}^{\infty}$ is \emph{tight} if for all $\varepsilon > 0$, there exists a compact set $K$ such that:
\begin{equation}
P(X_n \in K) > 1 - \varepsilon \quad \text{for all } n
\end{equation}
\end{definition}

\begin{theorem}[Tightness via Equicontinuity]
Let $\{f_n\}_{n=1}^{\infty}$ be a sequence of functions $f_n: \mathbb{R} \to \mathbb{R}$ satisfying:
\begin{enumerate}
    \item \textbf{Uniform Boundedness}: $\exists M > 0$ such that $|f_n(x)| \leq M$ for all $n, x$
    \item \textbf{Equicontinuity}: $\forall \varepsilon > 0, \exists \delta > 0$ such that $|x - y| < \delta \Rightarrow |f_n(x) - f_n(y)| < \varepsilon$ for all $n$
\end{enumerate}
Then $\{f_n\}$ is tight in the space of continuous functions.
\end{theorem}

\begin{proof}[Proof Sketch]
By the Arzelà-Ascoli theorem, boundedness and equicontinuity imply relative compactness in $C([0,1])$ with the uniform topology. For stochastic processes, this translates to tightness of the induced probability measures via Prokhorov's theorem.
\end{proof}

\subsection{Hölder Continuity}

\begin{theorem}[Hölder Continuity Implies Uniform Continuity]
Let $f: [0,1] \to \mathbb{R}$ be Hölder continuous with exponent $\gamma \in (0,1]$ and constant $C$:
\begin{equation}
|f(x) - f(y)| \leq C|x - y|^{\gamma}
\end{equation}
Then $f$ is uniformly continuous on $[0,1]$.
\end{theorem}

\begin{proof}
For any $\varepsilon > 0$, take $\delta = (\varepsilon/C)^{1/\gamma}$. Then $|x - y| < \delta$ implies:
\begin{equation}
|f(x) - f(y)| \leq C|x - y|^{\gamma} < C\delta^{\gamma} = \varepsilon
\end{equation}
\end{proof}

\section{Theorem 2: Spectral Decomposition}

\subsection{Mercer's Theorem}

\begin{theorem}[Mercer's Theorem]
Let $\Gamma: \mathbb{R} \times \mathbb{R} \to \mathbb{R}$ be a continuous, symmetric, positive semi-definite kernel. Then there exist:
\begin{enumerate}
    \item Eigenvalues $\{\lambda_m\}_{m=1}^{\infty}$ with $\lambda_m \geq 0$
    \item Orthonormal eigenfunctions $\{\phi_m\}_{m=1}^{\infty}$ satisfying $\int_{\mathbb{R}} \phi_m(x)\phi_n(x)\,dx = \delta_{mn}$
\end{enumerate}
such that:
\begin{equation}
\Gamma(s,t) = \sum_{m=1}^{\infty} \lambda_m \phi_m(s)\phi_m(t)
\end{equation}
where the series converges uniformly on compact sets.
\end{theorem}

\subsection{Weighted Chi-Square Distribution}

\begin{definition}[Weighted Chi-Square]
Let $\{\lambda_m\}_{m=1}^{\nu}$ be positive weights and $\{Z_m\}_{m=1}^{\nu}$ be i.i.d. standard normal variables. The weighted chi-square random variable is:
\begin{equation}
Q = \sum_{m=1}^{\nu} \lambda_m Z_m^2
\end{equation}
\end{definition}

\section{Theorem 3: Functional Delta Method}

\subsection{Hadamard Differentiability}

\begin{definition}[Hadamard Derivative]
Let $E, F$ be Banach spaces. A function $\phi: E \to F$ is \emph{Hadamard differentiable} at $x \in E$ with derivative $d\phi_x: E \to F$ if for all compact sets $K \subseteq E$:
\begin{equation}
\lim_{t \to 0} \sup_{h \in K} \frac{\|\phi(x + th) - \phi(x) - t \cdot d\phi_x(h)\|_F}{t} = 0
\end{equation}
\end{definition}

\begin{theorem}[Functional Delta Method]
Let $E, F$ be Banach spaces and $\{f_n\}_{n=1}^{\infty}$ a sequence of random elements in $E$ such that:
\begin{equation}
f_n \stackrel{d}{\longrightarrow} f \quad \text{as } n \to \infty
\end{equation}
If $\phi: E \to F$ is Hadamard differentiable at $f$, then:
\begin{equation}
\phi(f_n) \stackrel{d}{\longrightarrow} \phi(f) \quad \text{as } n \to \infty
\end{equation}
\end{theorem}

\section{Riemann Sum Approximation}

\subsection{Convergence of Double Sums}

\begin{theorem}[Riemann Sum Convergence]
Let $f: \mathbb{R}^2 \to \mathbb{R}$ be continuous on $[x_{\min}, x_{\max}] \times [y_{\min}, y_{\max}]$. For partitions with mesh size $\delta_n \to 0$:
\begin{equation}
\sum_{i,j} f(x_i, y_j) \Delta x_i \Delta y_j \longrightarrow \int_{x_{\min}}^{x_{\max}} \int_{y_{\min}}^{y_{\max}} f(x,y) \, dy \, dx
\end{equation}
\end{theorem}

\section{Conclusion}

The mathematical framework presented here provides the foundation for the Spatial CvM test formalization. Key results include:
\begin{itemize}
    \item Kernel smoothing preserves asymptotic properties under mixing conditions
    \item The empirical process converges weakly to a Gaussian process
    \item Mercer's decomposition connects kernel structure to chi-square limits
    \item The functional delta method enables inference on derived statistics
\end{itemize}

The Lean 4 formalization tracks these proofs, replacing axioms with constructive definitions where Mathlib provides adequate infrastructure, and documenting deep results from probability theory for future formalization.

\end{document}
